%%%%%%%%%%%%%%%%%%%%%%%%%%%%%%%%%%%%%%%%%%%%%%%%%%%%%%%%%%%%%%%%%%%%%%
\section{Outils logiciels envisagés}
%%%%%%%%%%%%%%%%%%%%%%%%%%%%%%%%%%%%%%%%%%%%%%%%%%%%%%%%%%%%%%%%%%%%%%

La première phase du projet ALORA repose essentiellement sur des outils
open source largement utilisés dans les domaines du développement
logiciel, de la gestion de données et du traitement automatique de la
parole.

Le choix de ces outils est motivé par leur robustesse, leur communauté
active ainsi que leur compatibilité avec les futures phases du projet.

%%%%%%%%%%%%%%%%%%%%%%%%%%%%%%%%%%%%%%%%%%%%%%%%%%%%%%%%%%%%%%%%%%%%%%

\subsection{Développement de la plateforme}

\begin{table}[H]

\centering

\begin{tabular}{p{4cm}p{9cm}}

\toprule

\textbf{Outil} & \textbf{Rôle} \\

\midrule

Flutter &
Développement de l'application de collecte (Android, iOS et Web). \\

FastAPI &
Développement de l'API REST assurant la communication entre
l'application et le serveur. \\

PostgreSQL &
Stockage des métadonnées relatives aux locuteurs, aux enregistrements et
aux annotations. \\

Git / GitHub &
Gestion des versions du projet et collaboration. \\

\bottomrule

\end{tabular}

\caption{Outils de développement}

\end{table}

%%%%%%%%%%%%%%%%%%%%%%%%%%%%%%%%%%%%%%%%%%%%%%%%%%%%%%%%%%%%%%%%%%%%%%

\subsection{Traitement des données}

\begin{table}[H]

\centering

\begin{tabular}{p{4cm}p{9cm}}

\toprule

\textbf{Outil} & \textbf{Rôle} \\

\midrule

Python &
Développement du pipeline de traitement des données. \\

FFmpeg &
Conversion, normalisation et vérification des fichiers audio. \\

Praat ou ELAN &
Transcription, segmentation et annotation des enregistrements. \\

Pandas &
Manipulation des métadonnées et préparation du corpus. \\

NumPy &
Traitement numérique des données. \\

Librosa &
Analyse et prétraitement des signaux audio. \\

\bottomrule

\end{tabular}

\caption{Outils de préparation du corpus}

\end{table}

%%%%%%%%%%%%%%%%%%%%%%%%%%%%%%%%%%%%%%%%%%%%%%%%%%%%%%%%%%%%%%%%%%%%%%

\subsection{Documentation}

La documentation technique du projet sera maintenue en parallèle du
développement.

Les principaux outils envisagés sont :

\begin{itemize}

\item LaTeX pour la rédaction des documents scientifiques ;

\item Markdown pour la documentation technique ;

\item GitHub pour le suivi des évolutions du projet.

\end{itemize}

%%%%%%%%%%%%%%%%%%%%%%%%%%%%%%%%%%%%%%%%%%%%%%%%%%%%%%%%%%%%%%%%%%%%%%

\subsection{Choix technologiques}

Les technologies retenues répondent aux objectifs suivants :

\begin{itemize}

\item privilégier des solutions open source ;

\item assurer une bonne évolutivité de la plateforme ;

\item faciliter l'intégration avec les outils de traitement automatique
de la parole ;

\item permettre la réutilisation du corpus par d'autres projets de
recherche.

\end{itemize}

Les choix présentés dans ce document constituent une première
proposition et pourront être adaptés en fonction des recommandations de
l'encadrant et des besoins identifiés au cours du développement.